\documentclass[a4paper,11pt,oneside,showtrims]{alpenthesis}
%%\aiiistock
\usepackage{lipsum}
\hexfalse
\paperfalse
\usepackage[numbered]{matlab-prettifier}

%% ================================================================= SET TITLE %
\title{My Thesis}
\author{Raphael Frey \\[1ex]\href{https://github.com/alpenwasser/}
                                 {\nolinkurl{https://github.com/alpenwasser/}}}

\chapterstyle{alpenthesis}
%% ============================================================== END PREAMBLE %
\begin{document}
%% ============================================================= BEGIN CONTENT %
%% <<< titlingpage
\begin{titlingpage}
    \fullhexpage{q1}{q0}
    \tikzset{external/export next=false}%
    \begin{tikzpicture}[remember picture,overlay,very thick]
        \node[draw,anchor=north east,yshift=-5mm,xshift=-10mm]
            at (current page.north east)
            {\texttt{LOGO TOP}};
    \end{tikzpicture}
    \tikzset{external/export next=false}%
    \begin{tikzpicture}[remember picture,overlay,very thick]
        \node[draw,anchor=south east,yshift=+5mm,xshift=-10mm]
            at (current page.south east)
            {\texttt{LOGO BOTTOM}};
    \end{tikzpicture}
    \flushright\sffamily

    \vspace{3ex}
    \Huge\bfseries{A Mostly Appropriate Title}\\[1ex]
    \Large\mdseries{Thesis}\\[3ex]

    \normalsize\mdseries

    alpenwasser\\
    team partner\\[3ex]

    Supervisor\\
    Expert\\

    \vfill

    \begin{tikzpicture}[very thick]
        \draw (0,0) rectangle (12,8);
        \draw (0,0) -- (12,8);
        \draw (0,8) -- (12,0);
        \node[anchor=center] at (6,1) {\LARGE\texttt{TITLE PICTURE}};
    \end{tikzpicture}

    \vfill

    \today\\
    Version 1.0.0
\end{titlingpage}
%% >>>
%% <<< frontmatter
\frontmatter
\tableofcontents*
\clearpage
\listoffigures*
\clearpage
\listoftables*
\clearpage
\listoflistings
%% >>>

\mainmatter
\part{Main Content}
\chapter{Text} % <<< Text
\lipsum[1]

\sffamily\lipsum[2]

\bfseries\lipsum[3]

\mdseries\ttfamily\lipsum[4]

\bfseries\ttfamily\lipsum[5]

\normalfont\scshape\lipsum[6]

\slshape\lipsum[7]

\scshape\lipsum[8]

\scslshape\lipsum[9]

\normalfont

%% >>>
\chapter{Ti\emph{k}Z Pictures, pgfplots} % <<< TikZ Pictures, pgfplots
\begin{figure}
    \centering
    \begin{tikzpicture}
        \coordinate (C) at (0,0);
        \coordinate (A) at (0,0);
        \coordinate (B) at (4,2);
        \draw (C) circle (1);
        \draw (A) rectangle (B);
    \end{tikzpicture}
    \caption{A Ti\emph{k}Z Picture}
\end{figure}

\begin{figure}
    \centering
    \begin{tikzpicture}
        \begin{axis}[width=\textwidth,height=50mm]
            \addplot[
                dv-7,
                domain=-3e-3:3e-3,
                samples=201,
            ]
            {exp(-x^2 / (2e-3^2)) / (1e-3 * sqrt(2*pi))};
        \end{axis}
    \end{tikzpicture}
    \caption{A plot of a function}
    \label{fig:function}
\end{figure}

\lipsum[3]
\begin{figure}[htb]
    \centering
    %\tikzexternaldisable
    \begin{tikzpicture}[trim axis left]
        \newcommand*\imp{IMPORTANT}
        % https://tex.stackexchange.com/a/78279/131649
        % https://tex.stackexchange.com/a/369884/131649
        % https://tex.stackexchange.com/q/36297/131649
        \pgfplotsset{
            every axis legend/.append style={
                at={(1,1.05)},
                anchor=south east,
            }
        }
        \begin{axis}[
                scale only axis,
                xtick = {
                    -0.0030,
                    -0.0015,
                     0,
                     0.0015,
                     0.0030
                },
                minor xtick = {
                    -0.0035,
                    -0.0025,
                    -0.002,
                    -0.001,
                    -0.0005,
                     0,
                     0.0005,
                     0.001,
                     0.002,
                     0.0025,
                     0.0035
                },
                extra x ticks={-0.0006}, % p.331ff in pgfplots manual
                extra x tick labels={\textcolor{dv+6}{\imp}},
                %xticklabel style={/pgf/number format/.cd,fixed,precision=5},
                xticklabel style={/pgf/number format/.cd,sci e}, % pgfplotstable: p.34
                extra x tick style={
                    grid=major,
                    tick label style={
                        rotate=90,
                        anchor=west,
                    },
                    tick pos = right,
                    ticklabel pos = right,
                },
                title style={
                    yshift={width("\imp")},
                },
                scaled ticks=false,
                width=\textwidth,
                height=80mm,
                xlabel=time (\si{\second}),
                ylabel=Amplitude (\si{\volt}),
                %xtick scale label code/.code={$\cdot \pi$},
                title=Plot Title,
                typeset ticklabels with strut,
                xlabel style={
                    yshift={width("\imp") - height("0")},
                },
                legend cell align=left,
                %legend pos = outer north east,
                every outer x axis line/.append style={black},
                every outer y axis line/.append style={black},
                axis line shift=0pt,
        ]

            \addplot[
                mark = *,
                ycomb,
                dv-7,
                domain=-3e-3:3e-3,
                samples=11,
                every mark/.style={fill=white},
            ]
                {exp(-x^2 / (2e-3^2)) / (1e-3 * sqrt(2*pi))};
            \label{pgfplots:bell}
            \addlegendentry{Bell Curve}

            \addplot[
                mark = *,
                ycomb,
                q0,
                samples=11,
                every mark/.style={fill=white},
                samples at={-0.0031,-0.0025,...,0.0029},
            ]
                {4.75e4 * x + 200};
            \label{pgfplots:linear}
            \addlegendentry{Some Linear Function}
        \end{axis}
    \end{tikzpicture}
    %\tikzexternalenable
    \caption{A plot similar to matlab's stem}
    \label{fig:stem}
\end{figure}

We have two curves in Figure~\ref{fig:stem}:
\begin{itemize}\tightlist
    \item[\ref{pgfplots:bell}] a bell curve,
    \item[\ref{pgfplots:linear}] and a linear function.
\end{itemize}
\lipsum[2]

%% >>>
\chapter{Tables} % <<< Tables
\begin{table}
    \centering
    \caption{tabular inside float}
    \label{tab:float}
    \begin{tabular}{lll}
        \toprule
        \scshape Header 1 & \scshape Header 2 & \scshape Header 3 \\
        \midrule
        Content           & Content           & Content           \\
        Content           & Content           & Content           \\
        Content           & Content           & Content           \\
        Content           & Content           & Content           \\
        \bottomrule
    \end{tabular}
\end{table}

\lipsum[3]

\begin{center}
    \tabcaption{Tabular outside of float}
    \label{tab:outside}
    \begin{tabular}{lll}
        \toprule
        \scshape Header 1 & \scshape Header 2 & \scshape Header 3 \\
        \midrule
        Content           & Content           & Content           \\
        Content           & Content           & Content           \\
        Content           & Content           & Content           \\
        Content           & Content           & Content           \\
        \bottomrule
    \end{tabular}
\end{center}
%% >>>
\chapter{References and Hyperlinks} % <<< References and Hyperlinks
This sentence refers to Table~\ref{tab:outside}.

This is a citation \cite{testitem}.

\href{https://hyperlink.com}{This is a hyperlink hiding behind text.}

\href{https://hyperlink.com}{\nolinkurl{https://hyperlink.com}}

%% >>>
\chapter{Sectional Headings} % <<<

This section illustrates the  style of \verb|\section|, \verb|\subsection| and
\verb|\subsubsection|.

\section{A Section}
\lipsum[2]

\subsection{A Subsection}
\lipsum[1]

\subsubsection{A subsubsection}
\lipsum[2]
%% >>>
\chapter{Boxes} % <<<
\tikzexternaldisable
\begin{tcolorbox}[breakable,title=Breakable Box With Title]
    \lipsum[1-9]
\end{tcolorbox}
\tikzexternalenable

\vspace{1em}

\tikzexternaldisable
\begin{commandshell}
    user:>  echo "test"
    root:> echo "test"
\end{commandshell}
\tikzexternalenable

\vspace{1em}

\tikzexternaldisable
\begin{tcolorbox}[title=Box with title]
    \lipsum[2]
\end{tcolorbox}
\tikzexternalenable

\vspace{1em}

\tikzexternaldisable
\begin{tcolorbox}
    \lipsum[2]
\end{tcolorbox}
\tikzexternalenable

\newpage
\tikzexternaldisable
Opening remarks \ldots
\begin{tcolorbox}[skin=octoboxfirst,title=Some Manually Broken Boxes]
    \lipsum[1]
\end{tcolorbox}
Intermittent remarks \ldots
\begin{tcolorbox}[skin=octoboxmiddle]
    \lipsum[2]
\end{tcolorbox}
More intermittent remarks \ldots
\begin{tcolorbox}[skin=octoboxlast]
    \lipsum[3]
\end{tcolorbox}
Closing remarks \ldots
\tikzexternalenable

%% >>>
\chapter{Code Listings} % <<<
\begin{minted}{tex}
    \begin{tabularx}{XX}
        a & a \\
        a & a \\
    \end{tabularx}
\end{minted}

\section{Matlab}
%%<<<

\begin{lstlisting}[
    style=Matlab-editor,
    gobble=4,
    title=Matlab Example Title,
    caption=Matlab Example Caption,
]
    %% Aufgabe 1: Knackgeraeusche entfernen (2/12)
    clear all; close all; clc;
    load Pruefung2_2.mat

    % ... %

    rxx = xcorr(x, scratch)/length(scratch);
    rxx = rxx((length(rxx)+1)/2:length(rxx));

    % c) Bestimmen Sie mittels plot() die Position der Knackgeraeusche. Welchen
    % Abstand weisen sie rund auf?
    figure; plot(rxx); grid on;
    % [px,py] = ginput(2);
    % px(2) - px(1)
    % ans = 1.4544e+04 = 14544

    % d) Legen sie aufgrund von c) einen Vektor p mit den Positionen der
    % Knackgeraeusche fest. Der Abstand darf als konstant angenommen werden.
    p = 14544:14544:length(x);

    figure; plot(x,'r');hold on;
    for i = p
        % i: scratch position
        src1 = x(i-512:i-257);
        src2 = x(i+256:i+511);
        x(i-256:i-1) = src1;
        x(i:i+255)   = src2;
    end
    plot(x,'b');

    %% Aufgabe 2: Distanzmessung (2/10)

    % ... %

    filter_koeffs = x; % filterkoeffizienten sind koeffizienten der stossantwort
    y0_upsampled = upsample(y0,10);
    y0_filtered  = filter(filter_koeffs,1,y0_upsampled);
    [peak0, offset] = max(y0_filtered);
    figure;plot(y0_filtered);grid on;

    % c) Bestimme Sie die Laufzeit und dann die Distanz zur Wand aufgrund von y1.
    y1_upsampled = upsample(y1,10);
    y1_filtered  = filter(filter_koeffs,1,y1_upsampled);
    [peak1, peak_index] = max(y1_filtered);
    peak_actual = peak_index - offset;
    figure;plot(y1_filtered);grid on;

    distance = T * peak_actual * c
\end{lstlisting}

%% >>>
%% >>>
\chapter{Mathematics} % <<<

A numbered equation:
\begin{equation}
    y(x) = x^2 + 2x + 5
\end{equation}

An unnumbered equation:
\begin{equation*}
    x_{01,02} = -1 \pm 2j
\end{equation*}

An \verb|align| with some numbered and unnumbered lines:

\begin{align}
    \frac{\Phi}{i}
    & =
    \int_{0}^{\infty} \frac{dx}{x} \int_{0}^{l+m}
    \left[
        \frac{y dy}{\sqrt{x^2+y^2}} - \frac{(y-l) dy}{\sqrt{x^2+(y-l)^2}}
    \right]
    \nonumber
    \\
    &=
    \int_{0}^{\infty}
    \left[
        \sqrt{x^2 + (l+m)^2)} - \sqrt{x^2+l^2} - \sqrt{x^2+m^2}+x
    \right]
    \frac{dx}{x}
    \nonumber
    \\
    &=
    \Bigg[
        \sqrt{x^2 + (l+m)^2} - \sqrt{x^2+l^2} - \sqrt{x^2+m^2}
        \nonumber
        \\
        & ~~~~~~~~~   + x - l \cdot \log{\frac{l+m+\sqrt{x^2+(l+m)^2}}{l+\sqrt{x^2+l^2}}}
        \nonumber
        \\
        & ~~~~~~~~~   - m \cdot \log{\frac{l+m+\sqrt{x^2+(l+m)^2}}{m+\sqrt{x^2+l^2}}} ~
    \Bigg]_0^{\infty}
    \label{eq:mISL:2}
    \\
    & \approx
    \left[ l \cdot \log{\frac{l+m}{l}} + m \cdot \log{\frac{l+m}{m}} \right]
    \label{eq:mISL:3}
\end{align}

\begin{equation}
    f(t) =
    \begin{cases}
        1 \quad t \geq 0\\
        0 \quad t < 0\\
    \end{cases}
\end{equation}

\begin{equation}
    % https://tex.stackexchange.com/q/369970/131649
    A \underset{f,g}{\times} B
\end{equation}

%% >>>
\cleardoublepage
%% <<< Appendices, Backmatter, Bibliography
\begin{titlingpage*}
    \fullhexpage{da2}{ct4}
    \begin{vplace}
        \flushright\Huge\bfseries\sffamily\appendixpagename
    \end{vplace}
\end{titlingpage*}
\appendix
\chapterstyle{alpenappendix}
\chapter{An Appendix Chapter}
\lipsum[1-3]

\chapter{Another Appendix Chapter with a Title That Requires Two Lines}
\lipsum[4-6]

\backmatter
\chapter{A Backmatter Chapter}
\lipsum[7-9]

\begin{thebibliography}{1}
    \bibitem{testitem}
        An Author, ``A Title``, 1979.
\end{thebibliography}
%% >>>
%% ============================================================== END CONTENT %
\end{document}
